% Gerado por regressao/06_tabelas_latex.R (projeto Modelagem). Não editar à mão.
\begin{landscape}
\begin{table}[htbp]
\caption{Robustez R11 — dependentes winsorizadas nos percentis 1 e 99 de cada ano}\label{tab:robustez_R11}
\centering
\footnotesize
\setlength{\tabcolsep}{3pt}
\begin{tabular}{@{}>{\raggedright\arraybackslash}p{\dimexpr\linewidth-594pt\relax}>{\centering\arraybackslash}p{60pt}>{\centering\arraybackslash}p{60pt}>{\centering\arraybackslash}p{60pt}>{\centering\arraybackslash}p{60pt}>{\centering\arraybackslash}p{60pt}>{\centering\arraybackslash}p{60pt}>{\centering\arraybackslash}p{60pt}>{\centering\arraybackslash}p{60pt}>{\centering\arraybackslash}p{60pt}@{}}
\hline
\textbf{Variável} & \textbf{\hspace{0pt}Despesa total} & \textbf{\hspace{0pt}Saúde e Saneamento} & \textbf{\hspace{0pt}Educação e Cultura} & \textbf{\hspace{0pt}Habitação e Urbanismo} & \textbf{\hspace{0pt}Assistência Social} & \textbf{\hspace{0pt}Transporte} & \textbf{\hspace{0pt}Administração} & \textbf{\hspace{0pt}Agricultura} & \textbf{\hspace{0pt}Comunicações} \\ \hline
\multicolumn{10}{@{}l}{\textit{Modelo A (ciclo eleitoral)}} \\
Ano eleitoral & 292,138\textsuperscript{***} & 57,927\textsuperscript{***} & 53,229\textsuperscript{*} & 127,166\textsuperscript{***} & 14,480\textsuperscript{***} & 40,899\textsuperscript{***} & 8,516 & 1,898 & $-$0,323\textsuperscript{***} \\
 & (94,415) & (9,612) & (28,734) & (29,799) & (4,100) & (9,144) & (9,160) & (1,940) & (0,072) \\
Ano pré-eleitoral & 155,066\textsuperscript{*} & 0,096 & 77,656\textsuperscript{**} & 31,279\textsuperscript{**} & 9,110\textsuperscript{*} & 6,739 & 11,063 & 5,251 & 0,712\textsuperscript{***} \\
 & (93,591) & (11,958) & (39,221) & (15,885) & (4,690) & (7,173) & (8,163) & (3,736) & (0,182) \\
\hline
Observações & 70.714 & 70.557 & 70.582 & 69.600 & 70.517 & 54.434 & 70.589 & 64.488 & 12.363 \\
Municípios & 5.568 & 5.568 & 5.568 & 5.564 & 5.568 & 5.205 & 5.568 & 5.436 & 1.897 \\
R\textsuperscript{2} & 0,830 & 0,757 & 0,717 & 0,336 & 0,403 & 0,241 & 0,346 & 0,199 & 0,047 \\[6pt]
\multicolumn{10}{@{}l}{\textit{Modelo B (coalizões)}} \\
Prefeito na coalizão & 35,886\textsuperscript{***} & 4,615\textsuperscript{**} & 0,173 & 6,484\textsuperscript{***} & 0,568 & 2,176 & 9,351\textsuperscript{***} & $-$0,702 & $-$0,184 \\
do governador & (6,680) & (2,311) & (3,024) & (2,514) & (0,680) & (2,090) & (2,372) & (0,939) & (0,237) \\
Prefeito na coalizão & $-$15,381 & 7,384\textsuperscript{**} & $-$2,117 & $-$0,481 & 0,679 & $-$10,389\textsuperscript{***} & $-$0,660 & $-$5,197\textsuperscript{***} & 0,532\textsuperscript{*} \\
do presidente & (10,238) & (3,321) & (3,998) & (3,477) & (0,950) & (2,585) & (3,392) & (1,272) & (0,313) \\
\hline
Observações & 70.714 & 70.557 & 70.582 & 69.600 & 70.517 & 54.434 & 70.589 & 64.488 & 12.363 \\
Municípios & 5.568 & 5.568 & 5.568 & 5.564 & 5.568 & 5.205 & 5.568 & 5.436 & 1.897 \\
R\textsuperscript{2} & 0,467 & 0,294 & 0,274 & 0,084 & 0,142 & 0,107 & 0,155 & 0,077 & 0,028 \\
\hline
\end{tabular}
\fonte{Elaboração própria.}
\nota{Modelo A: efeito fixo de município, erros-padrão de Driscoll-Kraay. Modelo B: efeitos fixos de município e de ano, erros-padrão agrupados por município. Erros-padrão entre parênteses. \textsuperscript{***}p$<$0,01; \textsuperscript{**}p$<$0,05; \textsuperscript{*}p$<$0,1.}
\end{table}
\end{landscape}
